\ssr{ЗАКЛЮЧЕНИЕ}

В ходе выполнения лабораторной работы №~7 была достигнута цель --- создана и протестирована программа на языке Python,
предназначенная для анализа текста, извлеченного из PDF-документов с использованием библиотеки PyPDF2 и регулярных выражений.
Были выполнены следующие задачи:
\begin{enumerate}
	\item разработаны регулярные выражения для решения задачи по варианту;
	\item реализована функция для поиска подстроки по варианту в pdf-файле с использованием разработанных регулярных выражений;
	\item реализовано программное обеспечение, принимающее на вход путь к pdf-файлу, использующее функцию из предыдущего пункта;
	\item проверена реализация на приложенных к лабораторной работе файлах;
	\item построена таблица с найдеными ошибками, каждой из реализаций.
\end{enumerate}

Код написанный, всеми моделями, допускал ложные срабатывания или не распозновал все ошибки. Наиболее точным оказался код предложенный моделью Qwen-Coder.
Код от модели DeepSeek находит больше всего ошибочных токенов, в то время как реализация от qwen срабатывает меньше всех. 
Таким образом, цель работы была достигнута, так как все поставленные задачи выполнены в полном объеме.